As the world’s largest athletics association, the Deutscher Leichtathletik-Verband (DLV) is tasked with facilitating peak individual performance and maintaining an internationally 
competitive national team \citep{DLV2026}. However, zero medals at the 2023 World Athletics Championships in Budapest and a historic low medal count at the 2024 Paris Olympics have left 
the organization in a sobering position. Although individual athletes still occasionally reach the podium, the declining success with regard to Olympic performances \citep{fremerey2024} raises questions about the fundamental strength of the national elite and extended elite.

While medals are the most visible metric of success, they reflect only the peak of elite performance and offer little insight into how performance is distributed within the elite level. 
To examine potential structural shifts, we analyze the top 35 rankings from 2001 to 2025 based on two research questions. First, we exploratively assess the impact of the COVID-19 pandemic as a short-term disruption to the performance development. Second, motivated by observed performance drop-offs beyond the top positions, we assess whether the performance gap 
within the elite has widened over time. Such a trend could indicate reduced internal competition, which can weaken international success, as elite development relies on strong 
competitive environments at the club level \citep[Tihanyi, 2001, cited in][]{Green2005}.

Our analytical approach is detailed in \Cref{sec:methods}, where we describe the data processing and justify the analytical framework. Building on this, we will analyze performance 
metrics to identify structural trends and evaluate our research questions in \Cref{sec:results}. Finally, \Cref{sec:conclusion} discusses the implications of these findings.